Un ejemplo claro de grandes p\'erdidas por riesgo de mercado es el caso de Controladora Comercial Mexicana (CCM) en 2008, cuando tom\'o posiciones en derivados cambiarios tipo \textit{Target Accrual Redemption Notes} (TARN). El problema fue que estos instrumentos parec\'ian convenientes mientras el tipo de cambio se manten\'ia estable, pero si el d\'olar superaba cierto nivel, la empresa quedaba obligada a vender al intermediario el doble del monto contratado a un precio desfavorable. Si el d\'olar permanec\'ia por debajo de 11.50 pesos, la empresa obten\'ia una ganancia; pero si rebasaba ese nivel, se activaba una cl\'ausula que multiplicaba la exposici\'on y pod\'ia generar p\'erdidas muy elevadas.

La falla principal en los controles fue que la operaci\'on dej\'o de ser una cobertura razonable y se convirti\'o en una posici\'on pr\'acticamente especulativa. Hubo una confianza excesiva en que el peso seguir\'ia estable y no se dimension\'o correctamente el tama\~no de la p\'erdida potencial. El propio material se\~nala que este caso reflej\'o una pr\'actica de administraci\'on de riesgos muy precaria y que la especulaci\'on empresarial incluso a\~nadi\'o presi\'on al tipo de cambio. Adem\'as, en las presentaciones tambi\'en se insiste en que los derivados complejos exigen controles muy estrictos porque son transacciones apalancadas y su valuaci\'on puede implicar un riesgo operativo considerable.

En cuanto a los protagonistas, el caso gir\'o alrededor de Comercial Mexicana, de su \'area de tesorer\'ia y finanzas, de los intermediarios financieros con los que celebr\'o los contratos y de las autoridades y especialistas que despu\'es analizaron el episodio. En el material aparecen, por ejemplo, las opiniones de Agust\'in Carstens, quien se\~nal\'o que la especulaci\'on de las empresas hab\'ia presionado adicionalmente al tipo de cambio, y de analistas como Francisco Su\'arez, quien calific\'o el caso como emblem\'atico por el tama\~no del quebranto y por la debilidad de la administraci\'on de riesgos.

La propuesta de soluci\'on ser\'ia establecer una pol\'itica mucho m\'as estricta de uso de derivados, permitiendo solo operaciones que realmente cubran una exposici\'on identificable de la empresa. Tambi\'en ser\'ia necesario fijar l\'imites de exposici\'on, aplicar pruebas de estr\'es para escenarios extremos de tipo de cambio y complementar el monitoreo con medidas como VaR y CVaR, ya que en tus diapositivas se se\~nala que la finalidad de las pruebas de estr\'es es detectar vulnerabilidades y asegurar que la entidad pueda soportar escenarios probables sin quebrar. Adem\'as, tendr\'ia que existir una separaci\'on clara entre quienes contratan los derivados y quienes supervisan el riesgo, para evitar que tesorer\'ia opere sin un control independiente.

Si yo hubiera sido el CRO, habr\'ia hecho algo distinto desde el inicio: no habr\'ia autorizado una estructura con p\'erdidas potencialmente desproporcionadas frente al beneficio esperado. Habr\'ia exigido que cada derivado estuviera vinculado a una necesidad real de cobertura, habr\'ia impuesto l\'imites de p\'erdida m\'axima aceptable y habr\'ia realizado escenarios de estr\'es considerando una depreciaci\'on fuerte del peso. En otras palabras, habr\'ia privilegiado una cobertura simple y transparente en lugar de aceptar un derivado ex\'otico que pod\'ia comprometer la estabilidad financiera de la empresa.